\documentclass[11pt]{article}
\usepackage[margin=1in]{geometry}
\usepackage{amsmath,amssymb,amsthm,mathtools}
\usepackage[hidelinks]{hyperref}

\newtheorem{theorem}{Theorem}
\newtheorem{lemma}{Lemma}
\newcommand{\Soc}{\operatorname{Soc}}
\newcommand{\rank}{\operatorname{rank}}

\title{Log-concavity of codimension-three level Hilbert functions of type two}
\author{}
\date{}

\begin{document}
\maketitle

\begin{abstract}
Let $R=k[x_1,x_2,x_3]$, where $k$ is any field.  We prove that the
Hilbert function
\[
  h=(1,3,h_2,\ldots,h_e=2)
\]
of every standard graded Artinian level quotient of $R$ of type two is
log-concave.  The argument uses the minimal resolution of the reindexed
graded Matlis dual.  Its last syzygy has rank one and full support; this
forces every proper subset of the middle columns to be independent over
$\operatorname{Frac}(R)$.  A putative failure of log-concavity then
reduces to a rank-one family of relations, which determines a cyclic
Gorenstein submodule.  Stanley's characteristic-free theorem for
codimension-three Gorenstein Hilbert functions closes the remaining case.
\end{abstract}

For $j\geq0$, set
\[
  N_j=\dim_k R_j=\binom{j+2}{2},
\]
and put $N_j=0$ for $j<0$.  We will use the following
characteristic-free form of Stanley's theorem.

\begin{lemma}[Stanley]
\label{lem:stanley}
If $B$ is a standard graded Artinian Gorenstein algebra of embedding
dimension at most three and socle degree $E$, then its Hilbert function
is symmetric and is nondecreasing through degree $\lfloor E/2\rfloor$.
\end{lemma}

More precisely, the Hilbert function is an SI-sequence; the
monotonicity stated above is the only part needed here.  The result
holds over an arbitrary field.  A characteristic-free proof is given
in F.~Zanello,
\emph{Stanley's theorem on codimension $3$ Gorenstein $h$-vectors},
Proc. Amer. Math. Soc. 134 (2006), 5--8,
\url{https://arxiv.org/abs/math/0411231}.

\begin{theorem}
Let $A=R/I$ be a standard graded Artinian level algebra of embedding
dimension three, socle degree $e$, and type two.  Then
\[
  h_i^2\geq h_{i-1}h_{i+1}
  \qquad(1\leq i\leq e-1).
\]
The statement holds in every characteristic.
\end{theorem}

\begin{proof}
Put $s=e+3$.  Since $A$ is level of type two with socle degree $e$,
the last free module in its minimal resolution is $R(-s)^2$.
Reindex the graded Matlis dual by setting
\[
  M_d=\operatorname{Hom}_k(A_{e-d},k),
  \qquad
  g_d=\dim_k M_d=h_{e-d}.
\]
Equivalently, $M\cong\operatorname{Ext}^3_R(A,R(-s))$.  Dualizing the
minimal resolution of $A$ gives a minimal exact complex
\begin{equation}
\label{eq:resolution}
  0\longrightarrow R(-s)
  \xrightarrow{\delta_3}F_2
  \xrightarrow{\delta_2}F_1
  \xrightarrow{\delta_1}R^2
  \longrightarrow M\longrightarrow0.
\end{equation}
Thus $M$ is generated by two elements in degree zero.  Its socle is
one-dimensional and concentrated in degree $e$.  Indeed, for $d<e$,
standard generation gives $A_{e-d}=R_1A_{e-d-1}$, so no nonzero
functional in $M_d$ is annihilated by all of $R_1$.

Write
\[
  F_1=\bigoplus_j R(-j)^{p_j},
  \qquad
  F_2=\bigoplus_j R(-j)^{q_j}.
\]
Let $K=\operatorname{Frac}(R)$.  Tensoring
\eqref{eq:resolution} with $K$ shows that the kernel of $\delta_2$ is
one-dimensional.  Its spanning vector has no zero coordinate:
$\delta_3$ is dual to the first differential in the resolution of
$A$, so its coordinates are precisely the minimal homogeneous
generators of $I$, and each is nonzero.
Consequently,
\begin{equation}
\label{eq:circuit}
  \text{every proper subset of the columns of $\delta_2$ is
  $K$-linearly independent.}
\end{equation}

Fix $d$ with $2\leq d\leq e$.  The corresponding log-concavity
inequality is the one centered at $g_{d-1}$.  Define
\[
  Q_d=\sum_{j\leq d}q_j,
  \qquad
  P_{<d}=\sum_{j<d}p_j,
\]
and let $r_d$ be the $K$-rank of the restriction of $\delta_1$ to the
summands of $F_1$ having shift strictly less than $d$.  Thus
$r_d\in\{0,1,2\}$.  Finally, set
\[
  \varepsilon_d=
  \begin{cases}
  1,&Q_d=\rank F_2,\\
  0,&Q_d<\rank F_2.
  \end{cases}
\]

By \eqref{eq:circuit}, the columns of $\delta_2$ having shift at most
$d$ have rank $Q_d-\varepsilon_d$.  Since every differential in a
minimal graded free resolution has entries in the homogeneous maximal
ideal, these columns can have nonzero entries only in rows of $F_1$
whose shifts are strictly less than $d$.  Their images lie in the
kernel of a rank-$r_d$ map from a $P_{<d}$-dimensional $K$-space to
$K^2$.  Therefore
\begin{equation}
\label{eq:rankbound}
  Q_d-\varepsilon_d\leq P_{<d}-r_d.
\end{equation}

Let
\[
  \Delta^2g_d=g_d-2g_{d-1}+g_{d-2}.
\]
The Hilbert numerator of \eqref{eq:resolution} is
\[
  (1-t)^3\sum_{j\geq0}g_jt^j
  =
  2-\sum_jp_jt^j+\sum_jq_jt^j-t^s.
\]
Since $d\leq e<s$, summing its coefficients through degree $d$ gives
\[
  \Delta^2g_d=2+Q_d-P_{<d}-p_d.
\]
Together with \eqref{eq:rankbound}, this yields the central estimate
\begin{equation}
\label{eq:central}
  \Delta^2g_d\leq
  2-r_d+\varepsilon_d-p_d.
\end{equation}

If $\Delta^2g_d\leq0$, then
$g_{d-2}+g_d\leq2g_{d-1}$, and the arithmetic--geometric mean
inequality gives
$g_{d-2}g_d\leq g_{d-1}^2$.  Thus a failure of log-concavity
would imply
\begin{equation}
\label{eq:positive}
  \Delta^2g_d\geq1.
\end{equation}

We now analyze the resulting cases.

First suppose $\varepsilon_d=1$.  Every summand $R(-b)$ of $F_2$
then has $b\leq d$.  By the degree correspondence in the dual
resolution, it comes from a minimal generator of $I$ of degree
$s-b\geq s-d$.  Thus $I$ has no nonzero element in degrees below
$s-d$.  With $m=e-d+2=s-d-1$, the tested triple is
\[
  (g_{d-2},g_{d-1},g_d)
  =(N_m,N_{m-1},N_{m-2}),
\]
Since the ratios $N_j/N_{j-1}$ are nonincreasing, this triple is
log-concave.

We may therefore assume $\varepsilon_d=0$.  If $r_d=2$,
\eqref{eq:central} gives $\Delta^2g_d\leq0$, contradicting
\eqref{eq:positive}.

If $r_d=0$, there are no summands of $F_1$ of shift below $d$.
Indeed, each such summand gives a nonzero minimal relation, so the
restricted map would have positive rank otherwise.  Equation
\eqref{eq:rankbound} then gives $Q_d=0$.  Consequently
\[
  g_{d-2}=2N_{d-2}=d(d-1),\qquad
  g_{d-1}=2N_{d-1}=d(d+1),
\]
and
\[
  g_d=2N_d-p_d=(d+1)(d+2)-p_d.
\]
A direct calculation gives
\[
  g_{d-1}^2-g_{d-2}g_d
  =
  d\bigl(2(d+1)+(d-1)p_d\bigr)>0.
\]

The only remaining case is $r_d=1$.  From
\eqref{eq:central} and \eqref{eq:positive}, a failure can occur only
if
\begin{equation}
\label{eq:hard}
  p_d=0,
  \qquad
  \Delta^2g_d=1.
\end{equation}

All columns of $\delta_1$ having shift below $d$ span one line over
$K$.  Because $R$ is a UFD, this line has a primitive homogeneous
generator $v=(v_1,v_2)\in R^2$, whose two components have a common
degree $a<d$.  The submodule generated by the low-shift columns is
then $vJ$ for a homogeneous ideal $J\subseteq R$.  To see this,
divide one nonzero column by the gcd of its two entries.  The resulting
primitive vector spans the line, and Euclid's lemma shows that every
other column is an $R$-multiple of it.

Because $p_d=0$, no relation has shift $d$, while relations of shift
greater than $d$ have no component in degrees at most $d$.  Therefore
\begin{equation}
\label{eq:kernel}
  \bigl(\ker(R^2\to M)\bigr)_t=(vJ)_t
  \qquad(t\leq d).
\end{equation}
Let $H$ be the image of $v$ in $M_a$ and put $n=d-a$.

If $H=0$, then $v$ itself is a relation of degree $a<d$.
Equation \eqref{eq:kernel} and the primitivity of $v$ force
$1\in J$, so $J=R$.  In this case set $B_j=0$ for every $j$.

Suppose $H\neq0$.  The finite-length module $RH$ has nonzero socle,
and this socle is contained in the one-dimensional socle of $M$.
Consequently
\[
  B=R/\operatorname{Ann}(H)
\]
is Artinian Gorenstein, has embedding dimension at most three, and has
socle degree
\[
  E=e-a.
\]
For every $j\leq n$ and $c\in R_j$,
\[
  cH=0
  \iff cv\in\ker(R^2\to M)_{a+j}
  \overset{\eqref{eq:kernel}}{\iff}cv\in(vJ)_{a+j}
  \iff c\in J_j.
\]
The last equivalence uses the injectivity of multiplication by the
nonzero vector $v$.  Thus $J$ and $\operatorname{Ann}(H)$ agree
through degree $n$.

In both the $H=0$ and $H\neq0$ cases, with negative subscripts
interpreted as zero, we have for $t=d-2,d-1,d$
\begin{equation}
\label{eq:formula}
  g_t=2N_t-N_{t-a}+B_{t-a}.
\end{equation}

Equation \eqref{eq:formula} gives
\[
  g_{d-1}\geq N_{d-1}.
\]
On the other hand,
\[
  g_{d-1}=h_{e-d+1}\leq N_{e-d+1}.
\]
Since $N_j$ is strictly increasing, it follows that
\begin{equation}
\label{eq:location}
  2d\leq e+2.
\end{equation}
When $H\neq0$, this gives
\[
  2(n-1)=2d-2a-2
  \leq e-2a
  \leq e-a=E.
\]
When $H\neq0$, Lemma~\ref{lem:stanley} gives
\[
  D:=B_{n-1}-B_{n-2}\geq0.
\]
For $H=0$ this also holds, since $D=0$ by definition.

Set
\[
  x=g_{d-1}-g_{d-2}.
\]
Using \eqref{eq:formula} and
$N_j-N_{j-1}=j+1$, we obtain
\[
  x=2d-n+D=d+a+D\geq d.
\]
Also
\[
  g_{d-2}\leq2N_{d-2}=d(d-1).
\]
Finally, \eqref{eq:hard} says that the tested triple has the form
\[
  (g_{d-2},g_{d-1},g_d)
  =(y,y+x,y+2x+1)
\]
for $y=g_{d-2}$.  Its log-concavity margin is
\[
  g_{d-1}^2-g_{d-2}g_d
  =(y+x)^2-y(y+2x+1)
  =x^2-y
  \geq d^2-d(d-1)
  =d>0.
\]
This contradicts the assumed failure.

We have proved log-concavity of $g$ at every index $d-1$ with
$2\leq d\leq e$.  Reversal preserves the log-concavity inequalities,
and $g_d=h_{e-d}$, so the same inequalities hold for $h$.

The argument is valid over every field.  The substantive theorem
invoked is the characteristic-free form of Stanley's theorem in
Lemma~\ref{lem:stanley}; the remaining steps use graded Matlis
duality, linear algebra over $\operatorname{Frac}(R)$, and the UFD
property of $k[x_1,x_2,x_3]$.
\end{proof}

\end{document}
